% ID: SEL-D0002-OA-E04
% ESTADO: cotejada. Fuente y transcripción contrastadas de forma independiente; salida revisada.
% Fuente: Selectividad/sel_2019_matematicas_aplicadas/Titular Septiembre/Exámenes Mat Apl CCSS II - A y B.pdf
% SHA-256: 31fb46947e141f5c8da8f7b381e7d27b67fe5a9d1597120454ce66fce0313f1a
% Página: 1; opción A; ejercicio 4.
% Encabezado original: EJERCICIO 4. Solo enunciado; no contiene solución.

Las puntuaciones obtenidas por los participantes en un concurso se distribuyen
siguiendo una ley Normal de varianza 36 y media desconocida. Se toma una
muestra aleatoria de 64 concursantes, cuya puntuación media es 35 puntos.

\begin{enumerate}[label=\alph*)]
  \item \textbf{(1.5 puntos)} Obtenga un intervalo, con un 92\,\% de confianza,
  para la puntuación media de los participantes en dicho concurso.
  \item \textbf{(1 punto)} Calcule el tamaño mínimo de la muestra que se ha
  de tomar para estimar la puntuación media del total de concursantes, con
  un error inferior a 2 puntos y un nivel de confianza del 98\,\%.
\end{enumerate}
